\chapter*{Resumen}
\addcontentsline{toc}{chapter}{Resumen}

Este trabajo desarrolla y evalúa un simulador modular de dispositivos
memristivos y componentes neuromórficos, junto con un agente de control
que utiliza un crossbar como representación computacional de sus pesos.
El sistema integra el modelo de Strukov, neuronas Leaky Integrate-and-Fire
(LIF), reglas de plasticidad, variabilidad estocástica y lectura de
crossbar mediante un solver nodal MNA. Con la definición residual $R^2=1-\mathrm{SSE}/\mathrm{SST}$ se obtienen
$R^2=0.999657$ para el estado interno y $0.996994$ para la corriente del
modelo de Strukov. Para STDP, un ajuste con nueve de los 13 puntos
digitalizados atribuidos a Jo et al. se evaluó en cuatro puntos
reservados y obtuvo $R^2=0.94220$ y RMSE $=2.136$ puntos porcentuales;
el tamaño reducido del holdout limita la evidencia de generalización.
La convención usada es $\Delta t=t_{\mathrm{post}}-t_{\mathrm{pre}}$,
positiva para pre antes que post y LTP; la definición en prosa del
artículo presenta el signo opuesto. Para LTP/LTD se simularon los pulsos
de $+3.2$ y $-2.8\,\mathrm{V}$, de $300\,\mu\mathrm{s}$, y lectura a
$1\,\mathrm{V}$ de la Fig.~2a. El test obtuvo $R^2=0.73988$,
RMSE $=8.618$ nA y ERP $=31.23\%$, con resultados deficientes por rama;
Pearson agregado es $r=0.99303$, aunque no implica concordancia de
amplitud y no está definido para la predicción constante de LTP. El
ajuste se hizo a la misma curva digitalizada y la evidencia es parcial.
Prezioso se presenta solo como reconstrucción visual aproximada,
sin métricas de ajuste ni asignación de ramas SET/RESET, porque la fuente
no especifica por completo el barrido temporal. Una curva
frecuencia-corriente LIF concuerda con la solución analítica del propio
modelo, mientras las pruebas D2D/C2C evalúan distribuciones sintéticas,
no mediciones de hardware. En un crossbar $2\times2$, el solver MNA
difiere 10.98\% de la VMM ideal al incluir resistencias de línea y
sensado. Los resultados respaldan la consistencia de varias
implementaciones numéricas, pero no una validación física general. En
CartPole, cinco corridas de 300 episodios con D2D/C2C habilitados
obtuvieron una recompensa media de 435.31 en los últimos 200 episodios
(rango: 369.66--474.31). La semilla 0 alcanzó tarde el nivel de
recompensa alta, por lo que su promedio de ventana final también refleja
el tiempo de aprendizaje; estos datos no establecen que una semilla sea
intrínsecamente mejor que otra. Tras seleccionar un checkpoint con 50
episodios de validación por matriz, una prueba con pesos congelados en 50 episodios
disjuntos obtuvo $493.63\pm0.93$ puntos. En Robot Sim, REINFORCE y R-STDP
usan el mismo camino simulado de crossbar, dispositivo vol\'atil de salida y
neuronas LIF. La pol\'itica de ambos se decodifica desde las corrientes de
columna; las LIF producen spikes y participan en la din\'amica WTA, pero el WTA
no restringe la acci\'on del R-STDP actual. Su elegibilidad usa la entrada
presin\'aptica y un evento de acci\'on centrado por la pol\'itica softmax; no
usa el spike LIF como evento postsin\'aptico. La regla modula la elegibilidad
con retornos futuros y actualiza el estado simulado $X$. Un piloto de 28
corridas compar\'o ambos algoritmos en CartPole, visi\'on y obst\'aculos, con
semillas 0 y 1 e inicializaciones $X_0=0.1$ y $0.9$. Se detuvo al obtener al
menos 271 exitos en la ventana m\'ovil de 300 episodios (90.33\%). R-STDP
alcanz\'o el criterio en 4/4 condiciones de cada tarea; REINFORCE, en 2/4 de
CartPole (dos registros incompletos), 4/4 de visi\'on y 4/4 de obst\'aculos.
Una repetici\'on de Car R-STDP, con evaluaci\'on cada 100 episodios, alcanz\'o
el criterio en 2/4 condiciones para semillas 2 y 4. Son resultados
exploratorios: hay dos semillas, las m\'etricas dependen del dominio y las
corridas originales de Car siguieron otra frecuencia de evaluaci\'on.

La plasticidad local R-STDP del crossbar simulado funciona como parte del
enfoque h\'ibrido que combina elegibilidad local con la se\~nal de recompensa
y asignaci\'on de cr\'edito de REINFORCE. En cambio, la variante que usa el
spike LIF real como evento post directamente en el aprendizaje de la tarea
tuvo rendimiento muy bajo en pruebas exploratorias. El piloto cuantitativo
corresponde al agente h\'ibrido etiquetado R-STDP, con elegibilidad
pre/acci\'on centrada por la pol\'itica; no a la variante de post LIF.
La fusi\'on act\'ua en la regla software de aprendizaje y recompensa; conserva
el mismo modelo y camino de lectura del crossbar.

\chapter*{Abstract}
\addcontentsline{toc}{chapter}{Abstract}

This work develops and evaluates a modular simulator for memristive
devices and neuromorphic components, together with a control agent that
uses a crossbar as a computational representation of its weights. The
system combines the Strukov model, Leaky Integrate-and-Fire (LIF)
neurons, plasticity rules, stochastic variability, and crossbar
readout through a nodal MNA solver. Evaluation combines digitized
references, analytical solutions, and reproducible numerical tests.
Using the residual definition $R^2=1-\mathrm{SSE}/\mathrm{SST}$, the
Strukov model obtains $R^2=0.999657$ for its internal state and
$0.996994$ for current. An exponential STDP rule was fitted to nine of
13 digitized points attributed to Jo et al. and evaluated on four
reserved points, yielding $R^2=0.94220$ and an RMSE of $2.136$ percentage
points; the small holdout limits the evidence for generalization. The
project uses $\Delta t=t_{\mathrm{post}}-t_{\mathrm{pre}}$ (positive for
pre-before-post and LTP); the paper's prose definition has the opposite
sign to its digitized positive-LTP axis. For Jo et al. (2010), LTP/LTD
was simulated with
the Fig.~2a pulse amplitudes ($+3.2$ and $-2.8\,\mathrm{V}$), widths
($300\,\mu\mathrm{s}$), and $1\,\mathrm{V}$ reads; the held-out result
was $R^2=0.73988$, RMSE $=8.618$ nA, and ERP $=31.23\%$, with poor
branch-wise scores and no independent-device prediction. The aggregate
Pearson coefficient is $r=0.99303$, but it does not imply amplitude
agreement; Pearson correlation is undefined for the constant LTP
predictions. The STDP comparison evaluates four reserved points.
Prezioso is reported only as
an approximate visual reconstruction, without fit metrics or branch
attribution, because the complete sweep protocol is unavailable. An LIF
frequency-current curve agrees with the analytical solution of the same
model, whereas D2D/C2C tests assess synthetic distributions rather than
hardware measurements. For a $2\times2$ crossbar, the MNA solver differs
by 10.98\% from the ideal VMM when line and sensing resistances are
included. These results support the numerical consistency of several
implementations, but do not establish general physical validation. In
CartPole, five 300-episode runs with D2D/C2C enabled achieved a mean
reward of 435.31 over their last 200 episodes (range: 369.66--474.31).
Seed 0 reached the high-reward regime later, so its final-window mean
also reflects learning time; these results do not establish that one
seed is intrinsically better than another. After selecting a checkpoint
using 50 validation episodes per matrix, a frozen-weight test on 50 disjoint episodes achieved $493.63\pm0.93$ points. In Robot Sim, both
REINFORCE and R-STDP use the same simulated crossbar, volatile output device,
and LIF pathway. Both policies decode actions from column currents; LIF
neurons generate spikes and participate in WTA dynamics, but WTA does not
gate the current R-STDP action selection. Its eligibility uses presynaptic
input and an action event centered by the softmax policy; it does not use the
LIF spike as its postsynaptic learning event. Reward-modulated eligibility
uses future returns to update the simulated state $X$. An exploratory 28-run
pilot compared both algorithms on CartPole, vision, and obstacle navigation
with seeds 0 and 1 and initial states $X_0=0.1$ and $0.9$. Runs stopped after
at least 271 successes in a rolling window of 300 episodes (90.33%). R-STDP
met the criterion in 4/4 conditions in each task; REINFORCE met it in 2/4
CartPole conditions (two records were incomplete), 4/4 vision conditions,
and 4/4 obstacle conditions. A separate Car R-STDP repeat, evaluated every
100 episodes, met the criterion in 2/4 conditions for seeds 2 and 4. These
are exploratory results: there are two seeds, success metrics are
domain-specific, and the original Car runs used a different evaluation
frequency.

The crossbar's local reward-modulated STDP works as part of the hybrid
approach, which combines local eligibility with REINFORCE's reward and
credit-assignment signal. By contrast, using the real LIF spike as the
postsynaptic event directly in task learning performed very poorly in
exploratory runs. The quantitative pilot corresponds to the hybrid agent
labelled R-STDP, using policy-centered pre/action eligibility, not the
real-LIF-post variant.
